\chapter{Pseudorandomness}\label{ch:oc-pseudorandom}

Volume~11 proved Roth's theorem by ergodic theory. Volume~14 proved it by a density increment. They are the same meeting: a set correlates with a character, or its $U^2$ norm is small. This chapter says that once, in the language of Chapters~\ref{ch:oc-integers} and~\ref{ch:oc-sieves}. There is no third proof of Roth. The series does not have nilsequences, and this chapter does not mention them as a tool it uses. Green--Tao stays in a heading marked as quoted.

\section{The $U^2$ count}

Volume~14, in its chapter on Gowers norms, proves the inverse theorem for the $U^2$ norm: for a function $f:\Z/N\Z\to\C$,
\[
\|f\|_{U^2}^4=\sum_{\xi}|\hat f(\xi)|^4,
\]
and if $|f|\le1$ then $\max|\hat f|\le\|f\|_{U^2}\le\max|\hat f|^{1/2}$. A set with small $U^2$ norm is a set whose balanced function has no large Fourier coefficient. The count of three-term progressions is then the random count.

\begin{theorem}\label{thm:u2count}
Let $A\subset\Z/N\Z$ with density $\delta$, and let $f=1_A-\delta$. Write
\[
\Lambda_3(A)=N^{-2}\#\{(x,d):x,\,x+d,\,x+2d\in A\}.
\]
Then
\[
\Lambda_3(A)=\delta^3+O\bigl(\|f\|_{U^2}\bigr).
\]
If $\|f\|_{U^2}$ is small, $A$ has asymptotically the number of three-term progressions of a random set of the same density. If it does not, Volume~14's inverse theorem supplies a character with which $f$ correlates, and the density-increment argument of that volume begins. This chapter does not run the increment.
\end{theorem}
\begin{proof}
Expand $1_A=\delta+f$. The average of $1_A(x)1_A(x+d)1_A(x+2d)$ splits into $\delta^3$ and seven terms each containing at least one $f$. A term with a single $f$ and two constants averages to $0$, because $\sum_x f(x)=0$. A term with two factors of $f$ is bounded by the $U^2$ norm through the identity of Volume~14: the Gowers inner product of $f$, $f$ and a bounded function is at most $\|f\|_{U^2}^2$, up to the normalisations of that chapter, and a term with three factors of $f$ is $\Lambda_3(f)$, which the same identity bounds by $\|f\|_{U^2}$. Collecting constants gives the $O(\|f\|_{U^2})$.
\end{proof}

The $W$-trick is the same idea as a sieve, read as a thinning. Volume~14 proves it: to count progressions in the primes one first passes to the integers not divisible by any prime $p<w$, by restricting to a residue class coprime to $W=\prod_{p<w}p$. In the language of Chapter~\ref{ch:oc-sieves} this is conditioning on the events $p\nmid n$ for $p<w$. The small primes are structure. After they are removed, the inverse theorem is applied to what remains. The trick is cited, not reproved.

\begin{examplex}[modulo $6$]
$W=6$. The classes coprime to $6$ are $1$ and $5$. A three-term progression of primes larger than $3$ lies in one of these classes, or jumps between them in a way the common difference records. Removing the primes $2$ and $3$ is the whole of the $W$-trick at this level.
\end{examplex}

The $U^2$ norm of an interval is large: an interval correlates with the trivial character. The $U^2$ norm of a random subset of density $\delta$ is small, of size about $N^{-1/4}$ times a power of $\delta$, by the fourth-moment calculation $\sum|\hat f|^4\asymp N^{-1}$ for a random $f$. A Bohr set, a set of the form $\{x:\|x\alpha\|<\varepsilon\}$, sits between them: it is the set where a character is biased, and Volume~14's inverse theorem says it is the model of every set with large $U^2$ norm.

\section{Green--Tao, quoted}

The Green--Tao theorem says that the primes contain arithmetic progressions of every finite length. The proof transfers Szemer\'edi's theorem from the integers to a pseudorandom majorant of the primes. Szemer\'edi's theorem and the transference are quoted from Volume~14, where they are already marked as quoted. The singular series for prime progressions, if Volume~14 numbered it, is cited from there; this chapter does not compute a new one.

What this chapter adds is the organisation. A deterministic set can match a mean and still be structured, which is Maier's theorem in Chapter~\ref{ch:oc-models}, or it can be pseudorandom in the $U^2$ norm and then have the random number of progressions, which is Theorem~\ref{thm:u2count}. The primes, after the $W$-trick, are conjectured to be the second kind for progressions of every length, and known to be so for length three by the two proofs in Volumes~11 and~14.
